\documentclass[11pt]{article}
\usepackage[margin=29mm]{geometry}
\usepackage[T1]{fontenc}
\usepackage{lmodern,microtype,amsmath,amssymb,amsthm,mathtools}
\usepackage[hidelinks]{hyperref}
\newtheorem{theorem}{Theorem}[section]
\newtheorem{lemma}[theorem]{Lemma}
\newtheorem{remark}[theorem]{Remark}
\newcommand{\Rea}{\operatorname{Re}}
\newcommand{\norm}[1]{\left\lVert#1\right\rVert}
\title{Near-edge extrapolation and parabolic smoothing\\
for finite-mode Navier--Stokes shear flows}
\author{C. K. Zeng\\
\small ORCID: \href{https://orcid.org/0009-0002-5402-3374}{0009-0002-5402-3374}}
\date{Scientific working draft: 5 September 2026}
\begin{document}
\maketitle

\noindent\textbf{Draft status.}
This is an incomplete, resumable scientific manuscript. The section below
records a complete supplementary lemma, independently checked at the
source-proof level. The main shear-flow results, their full proofs,
literature discussion, and bibliography have not yet been assembled here.
This file is not a frozen paper and is not approved for public release.

\section{A time-weighted estimate for decaying exponential sums}

The following elementary finite-dimensional estimate controls the onset
of a time cutoff without any separation assumption on the frequencies.
Its proof is included because a bound on the whole interval, in place of
the fixed initial interval used below, would introduce a spurious dependence
on its length. No claim of originality is made for this auxiliary lemma.

\begin{lemma}\label{lem:delay}
Let $(E,\mu)$ be a measure space, let $N\ge1$, and let
$\zeta_1,\ldots,\zeta_N\in\mathbb C$ be fixed with
$1\le\Rea\zeta_j\le4$. For measurable finite coefficients $c_j:E\to\mathbb C$,
set
\[
 f(t,z)=\sum_{j=1}^N c_j(z)e^{-\zeta_jt},\qquad
 k(t)=\int_E|f(t,z)|^3\,d\mu(z),\qquad
 K_T=\int_0^T k(t)\,dt.
\]
For every $T>0$ with $K_T<\infty$,
\[
 \int_0^T t\,k(t)^{2/3}\,dt\le C N^{4/3}K_T^{2/3},
\]
where $C$ is absolute. Coincident exponents are merged before the proof;
the displayed bound remains valid with the original $N$.
\end{lemma}

\begin{proof}
Write $a_j=\Rea\zeta_j$ and suppose first that the exponents are distinct.
The Laplace transform is an isometry from $L^2(0,\infty)$ onto the
right-half-plane Hardy space, with boundary norm
$(2\pi)^{-1}\int_{\mathbb R}|F(i\xi)|^2d\xi$.
In the span of $(s+\zeta_j)^{-1}$ an orthonormal basis is
\[
 \widehat\phi_j(s)=\frac{\sqrt{2a_j}}{s+\zeta_j}
       \prod_{r<j}\frac{s-\overline{\zeta_r}}{s+\zeta_r},
 \qquad 1\le j\le N.
\]
Here is the orthogonality argument. Each factor in the product has modulus
one on the imaginary axis and so is an isometric multiplier. After
cancelling the common initial factors in the inner product of the
$j$th and $\ell$th vectors, $\ell<j$, the remaining $j$th factor vanishes
at $\overline{\zeta_\ell}$. The reproducing identity for
$(s+\zeta_\ell)^{-1}$ makes the inner product zero. The last unmultiplied
kernel has squared norm $(2a_j)^{-1}$, giving unit normalization.
Partial fractions show that the first $j$ vectors lie in the first $j$
kernel span and that the coefficient of $(s+\zeta_j)^{-1}$ is nonzero:
the differences $\zeta_r-\zeta_j$ are nonzero, as are
$\zeta_j+\overline{\zeta_r}$ because their real parts are positive.
This also proves spanning.

For $V_\zeta h(t)=\int_0^t e^{-\zeta(t-u)}h(u)\,du$, inversion gives
\[
 \phi_j=\prod_{r<j}(I-2a_rV_{\zeta_r})
       \bigl(\sqrt{2a_j}e^{-\zeta_jt}\bigr).
\]
Expansion over subsets, $a_r\ge1$, and $2a_r\le8$ imply
\[
 |\phi_j(t)|\le\sqrt8 e^{-t}
       \sum_{\ell=0}^{j-1}\binom{j-1}{\ell}\frac{(8t)^\ell}{\ell!}
       =\sqrt8 e^{-t}L_{j-1}(-8t).
\]
Every $\ell$-fold convolution is bounded by
$e^{-t}t^\ell/\ell!$, which justifies the subset expansion even when the
$\zeta_r$ have different imaginary parts. For $x\ge0$,
\[
 L_j(-x)\le\sum_{\ell\ge0}\frac{(jx)^\ell}{(\ell!)^2}
 \le \cosh(2\sqrt{jx})\le e^{2\sqrt{jx}}.
\]
The middle inequality follows coefficientwise from
$\binom{2\ell}{\ell}\le4^\ell$.
For almost every $z$, let
$I_b(z)=\int_0^b|f(t,z)|^2dt$. Expansion in the orthonormal basis and
Cauchy--Schwarz give
\[
 |f(t,z)|^2
 \le8N e^{-2t+4\sqrt{8Nt}}I_\infty(z).
\]
Each $I_\infty(z)$ is finite because the coefficients are finite and all
real parts are positive.

Take $S=128N$. For $t\ge S$ the exponent is at most $-t$, whence
\[
 \int_S^\infty|f(t,z)|^2dt\le8Ne^{-S}I_\infty(z)
       <\tfrac12 I_\infty(z).
\]
Consequently $I_\infty\le2I_S$ and
$|f(t,z)|^2\le16Ne^{-t}I_S(z)$ for $t\ge S$.
If $T\ge S$, Minkowski in $L^{3/2}(E)$ and H\"older in time yield
\[
 \norm{I_S}_{L^{3/2}(E)}
 \le\int_0^S k(u)^{2/3}\,du
 \le S^{1/3}K_T^{2/3}.
\]
Thus
\[
 k(t)^{2/3}\le16NS^{1/3}e^{-t}K_T^{2/3}\quad(t\ge S).
\]
The contribution after $S$ is at most
$16NS^{1/3}(S+1)e^{-S}K_T^{2/3}\le C N^{4/3}K_T^{2/3}$.
Before $S$, H\"older directly gives
\[
 \int_0^{\min(T,S)} t\,k(t)^{2/3}dt
 \le \left(\int_0^{\min(T,S)}t^3dt\right)^{1/3}K_T^{2/3}
 \le4^{-1/3}S^{4/3}K_T^{2/3}.
\]
If $T<S$, this is already the entire assertion. Merging coincident
exponents only decreases the effective dimension.
\end{proof}

\begin{remark}[The cutoff application]
If the decay rates of $g(s,z)$ have real parts in $[\lambda,4\lambda]$,
$h(s)=\int_E|g(s,z)|^3d\mu$, $H=\int_0^4h(s)ds$, and
$0\le\chi(s)\le C_\chi s$, put $f(t,z)=g(t/\lambda,z)$.
Then $K_{4\lambda}=\lambda H$, and the lemma gives
\[
 \lambda\int_0^4\chi(s)h(s)^{2/3}ds
 \le C C_\chi N^{4/3}\lambda^{-1/3}H^{2/3}.
\]
This is a one-band temporal estimate, not a bound for arbitrary
Navier--Stokes solutions. It will be used in the full-plateau shear proof.
\end{remark}

\end{document}
